\section{Introduction}\label{sec:intro}

Large language models (LLMs) are increasingly used to produce not only snippets of code but complete, runnable artefacts: interactive scientific demonstrations \cite{chen2025interactscience}, physics animations \cite{wang2026simuscene,shende2026perceptual}, articulated 3D objects \cite{gao2026lam,le2025articulate}, CAD models \cite{alrashedy2025cad} and whole 3D worlds \cite{lu2026worldcoder}. One instance of this trend is the \emph{mechanism simulator}: a single web page in which a machine (a linkage, a gear train, a cam, a walking leg) is drawn in 3D and moves, with live numbers that explain it. Tools that generate such pages from a one-paragraph request are already available to end users \cite{teardown}. They are attractive for engineering education, design communication and, increasingly, as lightweight digital twins.

Such a page has a peculiar correctness problem. It almost always loads and animates, and it usually looks right. Whether it is \emph{kinematically} right---whether the parts stay connected, follow the motion the machine actually has, and respect the design rules that make the machine work---cannot be seen in a handful of frames and is not covered by any unit test, because there are none. In the generation pipelines we are aware of, the acceptance step is a vision-language model (VLM) that inspects screenshots \cite{shende2026perceptual,le2025articulate,alrashedy2025cad}. This is a mechanised version of the oldest simulation-validation technique, watching the animation \cite{sargent2013vv}, with the same weakness: plausible-looking motion is accepted.

Consider Theo Jansen's walking leg \cite{jansen_strandbeest}. Its eleven link lengths, the ``holy numbers'', make the foot trace a path with a long flat bottom. In our study, both models we tested remembered the eleven numbers but often assigned them to the wrong links. The resulting legs assemble, move smoothly and look like Jansen legs, but the flat stride shrinks from about a third of the cycle to a few percent. A VLM shown the frames reported the leg as correct in some runs and as faulty in others. When asked to re-check its own work, the generating model asserted that the lengths ``match the published holy numbers'' and, when told otherwise, that the checker must be wrong. An executable check that compares the declared lengths with the published ones settles the question in a fraction of a second and can also say which link is wrong and what its length should be.

This paper studies such executable oracles for LLM-generated mechanism simulators, both as an \emph{acceptance gate} and as a source of \emph{repair feedback}. The oracle (Section~\ref{sec:approach}) needs no hand-written specification per task. A minimal output contract asks the generator to expose a pure function \texttt{render(t)}, the dimensions it chose, and named marker points on each part. From these, the oracle derives metamorphic and structural checks (determinism under reordered calls, periodicity, rigid parts, coincident joints), checks the drawing against the declared dimensions, applies engineering design rules, and compares every marker with the kinematic solution computed from the declared dimensions. When a check fails, it reports where (which link, joint or marker, at which instant) and, where a ground truth exists, what the value should be.

We address four research questions:
\begin{itemize}
\item \textbf{RQ1} Which errors do LLMs make when generating mechanism simulators from free-form descriptions, and how often?
\item \textbf{RQ2} Which errors do the oracle and a VLM reviewer detect?
\item \textbf{RQ3} Does the oracle work as an acceptance gate for iterative repair, compared with VLM review in the loop and self-refinement?
\item \textbf{RQ4} How specific must the oracle's feedback be, and does the answer depend on the repairing model?
\end{itemize}

We study eight mechanisms, 88 generated simulators, 39 seeded mutants and nearly 300 repair attempts by two models of different size from the same family. The main findings are:
\begin{itemize}
\item Free-form specifications lead to errors in 20 of 80 simulators; detailed specifications lead to none in 8. Errors concentrate in design knowledge that must be recalled (Jansen's link lengths, Chebyshev's lambda proportions) and in assembly conditions. Mechanisms whose design rule can be stated as a relation (a cam's motion law, the Peaucellier--Lipkin conditions) produced no errors.
\item The oracle and the VLM catch different error classes. The oracle catches design-knowledge and invisible errors; the VLM catches geometry outside the oracle's contract. The VLM's detection of length errors depends largely on reading the numeric panel.
\item As a repair gate, oracle feedback that names the faulty element and its correct value fixes 28 of 28 faulty file-runs within two rounds. VLM review in the loop fixes 5 of 22 and self-refinement 7 of 28. VLM-driven loops stop early on wrong ``correct'' verdicts.
\item Specificity matters in a model-dependent way. On design-knowledge errors, the smaller model succeeded only with the correct value (13/13 versus 0/13 with location only). The larger model needed only the location (7/7) (Fisher's exact test, $p = 1.3\times10^{-5}$).
\end{itemize}

The contributions are: (i) an executable oracle for generated mechanism simulators that needs no per-task specification and includes a guard against marker gaming; (ii) a controlled comparison of oracle, VLM and self-refinement feedback in iterative repair with oracle-gated stopping; (iii) a feedback-specificity ablation showing an interaction between the information content of feedback and the repairing model; and (iv) the full dataset of specifications, generated and repaired simulators, mutants and oracle reports, available at \url{https://github.com/haihpse150218/llm-mechsim-oracle}.

Beyond education, the questions matter wherever simulation software is produced quickly and trusted on sight. This includes digital twins, operator training and design reviews, which are the kinds of systems in which NASA and other mission-critical organisations rely on simulation V\&V \cite{sargent2013vv,balci1997vva}. An acceptance gate that is cheap, deterministic and explains its verdict is a natural complement to human or model-based visual review.

Section~\ref{sec:background} gives background, Section~\ref{sec:approach} describes the oracle, Section~\ref{sec:design} the study design, Section~\ref{sec:results} the results, Section~\ref{sec:discussion} discusses them, Section~\ref{sec:threats} covers threats to validity, Section~\ref{sec:related} related work, and Section~\ref{sec:conclusion} concludes.
